\documentclass[11pt]{article}
\usepackage[margin=2.1cm]{geometry}
\usepackage{booktabs}
\usepackage{amsmath}
\usepackage{enumitem}
\usepackage{xcolor}
\usepackage[T1]{fontenc}
\usepackage{lmodern}

\setlength{\parindent}{0pt}
\setlength{\parskip}{3.6pt}
\newcommand{\code}[1]{\texttt{#1}}
\newcommand{\hr}{\vspace{2pt}\hrule\vspace{6pt}}

\begin{document}

\begin{center}
{\Large\bfseries EE\,362 --- Computer Organization Laboratory}\\[2pt]
{\large\bfseries Lab 6 --- Register File}\\[4pt]
Fall 2026 \quad$\cdot$\quad Lab TA: Gandhi Torres \quad$\cdot$\quad Sections: Mon / Tue, 3:45--5:45~PM, ECIII 237
\end{center}
\hr

\section*{1. Objective}

Build the processor's 32 general-purpose registers: two combinational read
ports, one clocked write port, and register 0 permanently zero.

\begin{quote}
\emph{Thirty-two copies of your Lab 2 register, plus address decoding.}
\end{quote}

\textbf{This is a lighter week than the ALU.} Do not be intimidated by ``32
registers''.

\section*{2. Background}

\subsection*{Three components you already own}

A register file is not a new kind of thing. It is:

\begin{center}
\begin{tabular}{lll}
\toprule
an address \textbf{decoder} & Lab 0 & which register does this address mean? \\
an array of \textbf{registers} & Lab 2 & each cell is your \code{RegN} \\
two \textbf{multiplexers} & Lab 3 & which register reaches each read port? \\
\bottomrule
\end{tabular}
\end{center}

Write those three equations down and the rest is assembly.

\subsection*{Specification --- Table C}

\begin{center}
\begin{tabular}{ccl@{\hskip 8pt}||@{\hskip 8pt}l}
\toprule
\# & \code{Strb} & \code{Dsel} & effect at the rising edge \\
\midrule
0 & 0 & --- & no write \\
1 & 1 & \code{"00000"} & \textbf{no write} --- register 0 is hardwired \\
2 & 1 & $k \neq 0$ & $R_k \leftarrow \code{Din}$ \\
\bottomrule
\end{tabular}
\end{center}

Reads are combinational and always valid:
$\code{Aout} = R[\code{Asel}]$, $\code{Bout} = R[\code{Bsel}]$, with
$R_0 \equiv 0$.

\subsection*{Equations}

\begin{align*}
\text{address minterm} \quad & m_k(S) = 1 \text{ iff } S = k
  && \text{(a 5-to-32 decoder)}\\
\text{write enable} \quad & \mathit{We}_k = \code{Strb} \cdot m_k(\code{Dsel}) \cdot [k \neq 0] \\
\text{storage} \quad & R_k^{+} = \mathit{We}_k \cdot \code{Din} + \overline{\mathit{We}_k} \cdot R_k
  && \text{(Lab 2, thirty-two times)}\\
\text{read port A} \quad & \code{Aout} = \sum_{k} m_k(\code{Asel}) \cdot R_k
  && \text{(a 32-to-1 mux)}
\end{align*}

\textbf{The $[k \neq 0]$ term is the entire mechanism of ``hardwired to zero''.}
Not a convention that software politely follows --- one extra AND condition, and
\code{\$0} becomes physically unwritable.

\subsection*{Why a permanent zero exists}

It makes the instruction set smaller. Several instructions come free:

\begin{center}
\begin{tabular}{ll}
\toprule
\code{move rd, rs} & is \code{ADD rd, rs, \$0} \\
\code{clear rd} & is \code{ADD rd, \$0, \$0} \\
\bottomrule
\end{tabular}
\end{center}

Your Lab 12 program relies on it: \code{SW r2, 0(r0)} computes its address as
$r_0 + 0 = 0$, with no special ``absolute address'' instruction needed.

\subsection*{Why two reads and one write}

Look at any R-type instruction: \code{ADD rd, rs, rt}. Two source registers
read, one destination written. The register file's shape is dictated by the
instruction format --- the ISA and the register file were designed against each
other.

Read ports are cheap: a read port is a multiplexer over storage that already
exists, and any number of readers can look at a register at once. Write ports
are not: two writers might target the same register in the same cycle, which
needs arbitration and a decision about what ``both wrote'' means. That is a new
correctness problem, not one more mux.

\subsection*{Read-during-write --- a decision you must be able to state}

If a read port addresses the register being written on the same edge, does it
see the old value or the new one?

\begin{center}
\begin{tabular}{ll}
\toprule
\textbf{READ-OLD} & what you get naturally: clocked write, combinational read \\
\textbf{READ-NEW} & needs an extra comparator and mux per read port (a bypass) \\
\bottomrule
\end{tabular}
\end{center}

Yours will be READ-OLD unless you deliberately build the bypass, and that is
correct for this CPU: in Lab 12 the register file is read at the start of a
cycle and written at the end of it.

\textbf{You must state which behaviour your design has in your report.}
``I do not know'' is not an acceptable answer about your own design.

\section*{3. Guided examples}

\begin{description}[leftmargin=2.4em,style=nextline,itemsep=2pt]
\item[Example 1 --- \code{regfile\_small}] $4 \times 8$, one read port. Small
enough to watch every register change in the waveform.
\item[Example 2 --- \code{regfile\_two\_ports}] The same file with a second read
port. The difference is two lines; compare the E4 numbers to see the cost.
\item[Example 3 --- \code{read\_during\_write}] Register 0, and the
read-during-write case, both asserted rather than assumed.
\end{description}

\section*{4. Assignment}

Complete \code{RegFile.vhd}: $32 \times 32$, ports \code{Asel}, \code{Bsel},
\code{Dsel}, \code{Din}, \code{Strb}, \code{clk}, \code{Aout}, \code{Bout}.

\textbf{Port names are fixed.} They come from the original EE 362 CPU library,
and the CPU in Lab 12 instantiates this entity by these exact names. Do not
rename them.

Block the write to register 0 \emph{and} force its read to zero. Either alone
would work; with both, a mistake in one place still cannot produce a non-zero
\code{\$0}.

\subsection*{Testbench}

Complete \code{tb\_Lab6.vhd}. The \code{wr} procedure makes a write one line.
\textbf{At least six checked vectors}, and your set must include all of:

\begin{itemize}[leftmargin=1.4em,itemsep=1pt]
\item write a register, read it back on \textbf{both} ports;
\item \textbf{try to write register 0}, then prove it still reads zero --- this
is the point of the lab;
\item a \textbf{read-during-write} case, checked \emph{before} and \emph{after}
the edge;
\item an edge with \code{Strb = '0'}, proving nothing is written;
\item a read that changes with \textbf{no clock edge}, proving reads are
combinational;
\item register 31, the boundary of the address range.
\end{itemize}

The supplied testbench already contains one failing vector. That is deliberate.

\section*{5. Evidence to collect (E1, E2, E4, E6, E7)}

\begin{center}
\begin{tabular}{lll}
\toprule
& Source & What it must show \\
\midrule
E1 & Flow Summary + Messages & 0 errors; warnings sorted \\
E2 & RTL Viewer & storage, and the read multiplexers \\
E4 & Resource Utilisation & \textbf{about 1024 registers} \\
E6 & Timing Analyzer & Fmax, and where the path runs \\
E7 & Waveform + transcript & register 0 refusing a write \\
\bottomrule
\end{tabular}
\end{center}

\textbf{This is the first design that SHOULD have registers.} Roughly 1024
($32 \times 32$). That number is a correctness check: far fewer usually means
the Fitter inferred block memory instead of flip-flops. Say which happened and
how you can tell.

\section*{6. Report --- \code{lab06.tex}}

\begin{itemize}[leftmargin=1.6em,itemsep=1pt]
\item \textbf{Specification:} Table C and the four equations.
\item \textbf{Design:} where the decoder, the registers and the multiplexers
are in your code; and \textbf{which read-during-write behaviour you built}.
\item \textbf{Results:} the register count, with a sentence on what it tells you.
\item \textbf{Discussion:} required. Must also answer: \emph{why is having a
permanently-zero register worth spending one of only 32 addresses on?}
\end{itemize}

\section*{7. Grading (100 points)}

\begin{center}
\begin{tabular}{lr}
\toprule
Pre-lab, submitted at the start of the session & 5 \\
Design compiles cleanly, warnings explained & 25 \\
Testbench and test vectors, coverage justified & 20 \\
Simulation results match the specification & 20 \\
Report chapter quality & 15 \\
In-lab demonstration and explanation & 15 \\
\midrule
Total & 100 \\
\bottomrule
\end{tabular}
\end{center}

Late work is accepted for one week at $-20$ points.

\section*{8. Things that will go wrong}

\begin{itemize}[leftmargin=1.6em,itemsep=1pt]
\item \textbf{Register 0 writable.} The most common miss, and it silently
corrupts Lab 12.
\item \textbf{Reads inside the clocked process.} Then reads become clocked, the
CPU cannot get its operands during a cycle, and Lab 12 fails in a way that is
very hard to see.
\item \textbf{Indexing with a vector.} \code{regs(Asel)} is meaningless ---
\code{to\_integer(unsigned(Asel))}.
\item \textbf{Off-by-one on the array range.} 32 registers is \code{0 to 31}.
\item \textbf{Renamed ports.} Lab 12 will not connect.
\end{itemize}

\vspace{10pt}
\hr
{\small This register file goes inside your CPU in Lab 12, addressed by the
instruction fields your decoder extracts in Labs 10--11, feeding the ALU you
built last week. Keep your files.}

\end{document}
